%=============================================================
% Résumé — Abstract — Remerciements
%=============================================================
\chapter*{Résumé}
\addcontentsline{toc}{chapter}{Résumé}
\thispagestyle{plain}

\begin{infobox}[Résumé exécutif]
Ce mémoire présente le projet MSPR \textit{Electio-Analytics}, réalisé dans le cadre du
titre RNCP 35584 (Bloc 3 — Big Data \& BI). L'objectif est de prédire le comportement
électoral de la région \textbf{Nouvelle-Aquitaine} (12 départements, 627 cantons) lors de
l'élection présidentielle 2022, à partir d'indicateurs socio-économiques issues des
recensements INSEE.
\end{infobox}

\vspace{0.5cm}

Le projet mobilise un pipeline de traitement de données en \textbf{8 phases} centré
sur \textbf{MongoDB} : collecte via internet (wget/Python), ingestion dans
\texttt{raw\_data}, nettoyage et ingénierie des variables ($\Delta$) depuis MongoDB,
stockage de la base finale dans \texttt{data\_final}, pipeline ML qui lit depuis
\texttt{data\_final} et stocke prédictions et graphiques dans MongoDB, et restitution
via une interface React 19 connectée à MongoDB.

Le jeu de données final comprend \textbf{40 000 lignes $\times$ 171 colonnes}, avec une
variable cible à 5 classes synthétiques (\textit{Crise, Déclin, Stable, Croissance, Boom})
reflétant la dynamique socio-économique locale.

Le modèle \textbf{HistGradientBoosting} obtient la meilleure performance avec une accuracy
de \textbf{82,11\%} et un F1-score macro de \textbf{82,10\%}. Appliqué aux 12 départements,
il prédit correctement le vainqueur dans \textbf{11 cas sur 12} (91,7\% de précision
départementale), avec une probabilité régionale d'\textbf{Emmanuel Macron à 68,7\%}.

\vspace{0.5cm}
\noindent\textbf{Mots-clés :} Machine Learning, Prédiction électorale, Nouvelle-Aquitaine,
HistGradientBoosting, Big Data, MongoDB, Pipeline ETL, React 19.

\bigskip
\hrule
\bigskip

\chapter*{Abstract}
\addcontentsline{toc}{chapter}{Abstract}
\thispagestyle{plain}

This report presents the MSPR \textit{Electio-Analytics} project, developed as part of
the RNCP 35584 certification (Block 3 — Big Data \& BI). The goal is to predict the
electoral behavior of the \textbf{Nouvelle-Aquitaine} region (12 departments, 627 cantons)
in the 2022 French presidential election, using socio-economic indicators from INSEE census data.

The project implements an \textbf{8-phase data processing pipeline centred on MongoDB}:
internet collection via wget/Python, raw ingestion into \texttt{raw\_data}, cleaning and
delta feature engineering reading from MongoDB, final storage in \texttt{data\_final},
ML pipeline loading from \texttt{data\_final} and storing predictions and charts back in
MongoDB, and delivery through a React 19 dashboard connected to MongoDB.

The final dataset contains \textbf{40,000 rows $\times$ 171 columns}, with a 5-class synthetic
target variable (\textit{Crise, Déclin, Stable, Croissance, Boom}) reflecting local
socio-economic dynamics.

The \textbf{HistGradientBoosting} model achieves the best performance with \textbf{82.11\%}
accuracy and \textbf{82.10\%} macro F1-score. Applied to the 12 departments, it correctly
predicts the winner in \textbf{11 out of 12 cases} (91.7\% departmental accuracy), with
a regional probability of \textbf{Emmanuel Macron at 68.7\%}.

\vspace{0.5cm}
\noindent\textbf{Keywords:} Machine Learning, Electoral Prediction, Nouvelle-Aquitaine,
HistGradientBoosting, Big Data, MongoDB, ETL Pipeline, React 19.

\bigskip
\hrule
\bigskip

\chapter*{Remerciements}
\addcontentsline{toc}{chapter}{Remerciements}
\thispagestyle{plain}

Nous tenons à remercier chaleureusement l'ensemble de l'équipe pédagogique d'EPSI
Bordeaux pour l'encadrement et les conseils prodigués tout au long de ce projet MSPR.

Nos remerciements vont également à l'équipe projet pour l'implication collective,
la rigueur analytique et la créativité déployées pour mener à bien cette mission.

Ce projet nous a permis de consolider des compétences transversales en ingénierie
des données, en machine learning et en développement d'interfaces interactives,
dans un contexte applicatif concret et socialement pertinent.
